\chapter{Marca, moedas e portas de compra e venda}
\noindent\textit{Vinheta.} Um pequeno comércio do bairro passa a aceitar apenas pagamentos digitais por ``segurança''. Na semana seguinte, o banco atualiza os termos: ``\emph{Para sua proteção, certas categorias terão limites automáticos}''. Um mês depois, uma política pública oferece \emph{descontos programáveis} para quem compra em parceiros credenciados. Conveniência? Sim. Mas uma pergunta NIC acende: \textbf{quem controla a porta}?

\section{Mapa bíblico}
\begin{itemize}[leftmargin=*,noitemsep]
  \item \textbf{Portas econômicas} (Ap~13:16--17): sem o sinal, ``ninguém compre ou venda'' --- \emph{culto} e \emph{coerção} se tocam.
  \item \textbf{Justiça e pesos honestos} (Pv~11:1; Mq~6:11): Deus se importa com as \emph{regras do mercado}.
  \item \textbf{Generosidade na escassez} (At~11:27--30; 2~Co~8): a igreja responde ao controle com \emph{partilha e fidelidade}.
\end{itemize}

\section{CBDCs e dinheiro programável (em linguagem simples)}
\textbf{CBDC} é moeda digital emitida pelo \emph{banco central}. Pode ser \textbf{de atacado} (entre bancos) ou \textbf{de varejo} (para o público). ``Programável'' significa que \emph{regras de uso} podem ser embutidas no pagamento (p.\,ex., validade, destinos permitidos, descontos condicionais). Benefícios prometidos: eficiência, inclusão, combate a fraude. Riscos: \textbf{portas} de compra e venda \emph{centralmente} controláveis, com erros e abusos possíveis.

\section{Panorama 2025 (resumo, sem alarmismo)}
\begin{itemize}[leftmargin=*,noitemsep]
  \item \textbf{Lançados} (uso doméstico em curso): Bahamas (Sand Dollar), Jamaica (Jam-Dex), Nigéria (eNaira).
  \item \textbf{Pilotos relevantes}: China (e-CNY, maior piloto), Índia (e-rupee, varejo em piloto), Hong Kong (e-HKD, piloto em fases), projetos de atacado para \emph{pagamentos transfronteiriços} (ex.: mBridge). 
  \item \textbf{Europa}: debate legislativo sobre \emph{euro digital}; Suécia concluiu pilotos técnicos do \emph{e-krona} e segue avaliando.
  \item \textbf{Brasil}: \textbf{Drex} --- arquitetura com \emph{camada de atacado} (liquidação em dinheiro do BC) e \emph{depósitos tokenizados} para o varejo (não é ``carteira do BC'' para o público). Foco: tokenização, eficiência e contratos inteligentes.
\end{itemize}

\begin{nicbox}
\textbf{NIC aplicado}\\
\textbf{Narrativa}: ``mais segurança e inclusão'' --- verdadeiro em partes, mas incompleto se omitir riscos de governança.\\
\textbf{Identidade}: carteira \emph{KYC} e reputações podem redefinir quem ``pode'' transacionar.\\
\textbf{Controle}: portas de compra/venda podem ser abertas/fechadas por políticas, algoritmos e \emph{parceiros} de plataforma.
\end{nicbox}

\section{Lei 2 do Engano Digital}
\textbf{Quem controla a porta, controla o culto}. Quando acesso econômico depende de \emph{conformidade comportamental}, a “porta” vira \textbf{altar}. O perigo não é a \emph{tecnologia} em si, mas a \textbf{lealdade} que ela exige para funcionar.

\section{Cenários práticos (o bom, o difícil e o perigoso)}
\begin{itemize}[leftmargin=*,noitemsep]
  \item \textbf{O bom}: benefícios pagos instantaneamente a quem mais precisa; \emph{cashback} para merenda escolar; taxas menores para microcomércios.
  \item \textbf{O difícil}: \emph{falsos negativos} e bloqueios automáticos deixam inocentes sem acesso por dias.
  \item \textbf{O perigoso}: \emph{restrições geográficas} e \emph{categorias proibidas} (ou “desincentivadas”) definidas de modo opaco; exclusão silenciosa de dissidentes ou vulneráveis.
\end{itemize}

\section{Brasil em foco (Drex, Pix e tokenização)}
\begin{itemize}[leftmargin=*,noitemsep]
  \item \textbf{Pix} segue como infraestrutura-chave para pagamentos instantâneos; tende a integrar-se com arranjos mais complexos (ex.: liquidação de ativos tokenizados).
  \item \textbf{Drex} propõe uma \emph{camada de atacado} com dinheiro do BC para liquidação e uma \emph{camada de varejo} por depósitos tokenizados via bancos --- com \textbf{contratos inteligentes} para programar regras \emph{no ativo}, não \emph{no cidadão}.
  \item \textbf{Implicação pastoral}: excelente para baratear \emph{boas} regras (fraude, conciliação); risco se “regras de vida” forem embutidas no dinheiro (uso, destino, prazo).
\end{itemize}

\section{Exemplos internacionais (instantâneo 2025)}
\begin{itemize}[leftmargin=*,noitemsep]
  \item \textbf{China (e-CNY)}: piloto em grande escala, incluindo transporte, varejo e eventos; discussão sobre uso transfronteiriço.
  \item \textbf{Índia (e-rupee)}: piloto de varejo com expansão para bancos e \emph{fintechs}; planos de explorar casos transfronteiriços.
  \item \textbf{Hong Kong (e-HKD)}: piloto por fases (\emph{e-HKD+}) com testes de atacado e varejo integrados.
  \item \textbf{mBridge (transfronteiriço)}: plataforma multi-BCs (China, Hong Kong, Tailândia, EAU; adesões adicionais) em estágio de produto mínimo viável para pagamentos internacionais.
  \item \textbf{Europa}: \emph{euro digital} em discussão política e desenho técnico; foco em privacidade e limites de saldo.
\end{itemize}

\section{Ferramenta --- ``Portas'' do seu orçamento (auditoria de 30 minutos)}
\begin{checklist}
\begin{itemize}[leftmargin=*,noitemsep]
  \item \textbf{Multi-rails}: tenha ao menos \emph{3 vias} de pagamento (Pix/cartão/dinheiro). Dinheiro vivo para emergências.
  \item \textbf{Reserva de liquidez}: 2--4 semanas de despesas essenciais em formas \emph{acessíveis offline}.
  \item \textbf{Documentos e prova de vida}: guarde alternativas \emph{não biométricas} para validações (cartório, senhas de recuperação).
  \item \textbf{Comércio local}: incentive negócios do bairro a manter ao menos um meio de pagamento \emph{fora de nuvem}.
  \item \textbf{Caridade organizada}: fundo da igreja para socorrer bloqueios temporários (alimentos, remédios, transporte).
\end{itemize}
\end{checklist}

\section{Ferramenta --- Política de generosidade e objeção de consciência (igreja e famílias)}
\begin{checklist}[title={Princípios práticos}]
\begin{itemize}[leftmargin=*,noitemsep]
  \item \textbf{Transparência}: explique por que usam certos meios de pagamento e quais dados são coletados.
  \item \textbf{Alternativa digna}: sempre ofereça \emph{um caminho} sem app ou biometria.
  \item \textbf{Limites}: recuse \emph{condições morais} embutidas no dinheiro; priorize \emph{servir pessoas} acima de metas algorítmicas.
  \item \textbf{Resiliência comunitária}: crie uma rede de apoio para membros que sofrerem \emph{bloqueios} inesperados.
\end{itemize}
\end{checklist}

\section{Perguntas para grupos pequenos}
\begin{itemize}[leftmargin=*,noitemsep]
  \item Onde sua família está \textbf{superdependente} de uma única ``porta'' de pagamento?
  \item Que \textbf{política de alternativa analógica} sua igreja precisa escrever esta semana?
  \item Em que ponto \emph{conveniência} começou a exigir \emph{lealdade} no seu coração?
\end{itemize}

\bigskip
\noindent\textit{Próximo capítulo: Transumanismo --- ``sereis como Deus?'' bio/neurotecnologia e limites éticos.}
